\section{Methodology}

\subsection{Problem Formulation}

Let

\[
K = \{k_1, k_2, \ldots, k_n\}
\]

denote the set of capabilities available to a system.

A capability configuration is a subset

\[
C \subseteq K.
\]

Executing configuration \(C\) produces a stochastic binary outcome

\[
Y(C) \in \{0,1\},
\]

where \(Y(C)=1\) denotes a failure.

Each configuration therefore has an unknown failure probability

\[
p(C)=P(Y(C)=1).
\]

The discovery problem is to identify minimal subsets of capabilities whose joint activation causes a substantial increase in failure risk, while minimizing the total number of stochastic executions required.

\subsection{Stochastic Capability Interactions}

A capability interaction is not defined only by a high absolute failure probability.

A candidate set must also exhibit risk beyond that observed in its relevant lower-order subsets.

For a candidate interaction \(C\), SCIF uses a joint-risk increment concept of the form

\[
JRI(C)
=
\hat{p}(C)
-
\max_{S \subset C}
\hat{p}(S),
\]

where \(\hat{p}\) denotes an empirical failure-rate estimate.

A candidate is therefore supported when its joint failure rate is sufficiently large and its observed risk cannot be explained by a lower-order subset.

\subsection{BSIB Benchmark Model}

The Benchmark for Stochastic Interaction Bugs (BSIB) provides synthetic capability configurations with hidden interaction faults.

The simulator exposes only stochastic execution outcomes to the discovery algorithm.

The underlying true failure probability is not included in the execution result.

This prevents the discovery algorithm from directly accessing the benchmark ground truth.

\subsection{Keyed Stochastic Simulation}

Experiments use deterministic keyed random streams.

For a fixed benchmark configuration and experiment seed, the random outcome stream is stable even when configurations are evaluated in a different order.

This design prevents algorithmic execution order from unintentionally changing the stochastic evidence observed for the same configuration.

The keyed simulator therefore supports reproducible and fair comparisons among discovery methods.

\subsection{Pairwise Discovery}

SCIF v0.0.4 begins with the SCIF v0.0.3 pairwise discovery procedure.

The frozen evaluation configuration uses:

\begin{itemize}
\tightlist
\item
  100 initial screening trials;
\item
  300 screening-retest trials;
\item
  1,500 confirmation trials;
\item
  1,000 subset-confirmation trials;
\item
  minimum joint failure rate of 0.15;
\item
  minimum joint-risk increment of 0.10; and
\item
  confidence threshold of 0.95.
\end{itemize}

Initial sampling provides a low-cost estimate of the failure signal.

Pairs with sufficiently convincing evidence proceed to confirmation.

Borderline configurations may receive additional screening trials before a final decision is made.

This architecture concentrates stochastic executions on configurations that show evidence of interaction risk.

\subsection{Limitation of Pairwise-Only Discovery}

A pure higher-order interaction can remain invisible to every lower-order subset.

For example, consider a three-way interaction involving capabilities

\[
\{a,b,c\}.
\]

It is possible for

\[
p(\{a,b\}),
p(\{a,c\}),
p(\{b,c\})
\]

to remain near baseline while

\[
p(\{a,b,c\})
\]

is substantially elevated.

A pairwise-only algorithm therefore cannot represent or recover this type of fault.

This limitation motivates the higher-order stages of SCIF v0.0.4.

\subsection{Known-Pair Suppression}

After pairwise discovery, SCIF v0.0.4 constructs probe configurations intended to disable already-discovered pairwise interactions.

For the set of discovered pairwise interactions, the algorithm computes minimal capability-removal sets that intersect every known pair.

Removing such a set creates a configuration in which the discovered pairwise interactions are suppressed.

Multiple minimal removal sets may exist.

SCIF evaluates the resulting residual configurations rather than assuming that one particular suppression is sufficient.

\subsection{Residual-Risk Detection}

Let \(C_r\) denote a configuration after suppression of known pairwise interactions.

SCIF estimates

\[
\hat{p}(C_r)
\]

and compares the remaining failure signal with the estimated baseline.

Higher-order escalation is triggered only when residual failure risk remains sufficiently large.

The frozen parameters are:

\begin{itemize}
\tightlist
\item
  1,000 residual trials;
\item
  minimum residual failure rate of 0.15; and
\item
  minimum residual-risk increment of 0.10.
\end{itemize}

Conceptually, the residual stage asks:

\begin{quote}
After explaining and suppressing the interactions already found, does substantial unexplained risk remain?
\end{quote}

If the answer is no, SCIF terminates without invoking higher-order localization.

If the answer is yes, the residual configuration becomes the input to the higher-order localizer.

\subsection{Residual Higher-Order Localization}

Higher-order localization is applied only after residual-risk detection has justified escalation.

For each capability in the selected residual configuration, SCIF measures the failure rate after removing that capability.

Let \(C_r\) be the source configuration and let \(k \in C_r\).

The empirical removal drop is

\[
D(k)
=
\hat{p}(C_r)
-
\hat{p}(C_r \setminus \{k\}).
\]

Capabilities whose removal causes a sufficiently large decrease in failure risk are treated as candidate members of the hidden interaction.

The frozen localization configuration uses:

\begin{itemize}
\tightlist
\item
  1,000 trials per configuration;
\item
  minimum failure rate of 0.15;
\item
  minimum removal drop of 0.10; and
\item
  minimum higher-order candidate size of three.
\end{itemize}

\subsection{Minimality Confirmation}

Removal evidence alone is not sufficient.

The resulting candidate is evaluated directly, and its immediate subsets are also tested.

The purpose is to verify that the candidate itself retains elevated failure risk while dropping one of its essential members removes the interaction signal.

This step protects against reporting unnecessarily large interaction sets.

\subsection{Complete SCIF v0.0.4 Pipeline}

The final procedure is:

\begin{enumerate}
\def\labelenumi{\arabic{enumi}.}
\tightlist
\item
  perform adaptive pairwise discovery;
\item
  record confirmed pairwise interactions;
\item
  construct minimal suppression sets for the known pairs;
\item
  evaluate residual-risk probes;
\item
  stop when no substantial residual risk remains;
\item
  otherwise select a residual configuration;
\item
  apply deletion-style higher-order localization;
\item
  confirm candidate minimality; and
\item
  return pairwise and higher-order discoveries in a unified report.
\end{enumerate}

The higher-order stage is therefore conditional rather than mandatory.

This design preserves the efficiency of pairwise discovery on simple scenarios while providing an escalation path for unexplained higher-order risk.
